Дано расстояние $L$ в сантиметрах. Используя операцию деления нацело, найти количество полных метров в нем (1 метр = 100 см).
Дана масса $M$ в килограммах. Используя операцию деления нацело, найти количество полных тонн в ней (1 тонна = 1000 кг).
Дан размер файла в байтах. Используя операцию деления нацело, найти количество полных килобайтов, которые занимает данный файл (1 килобайт = 1024 байта).
Даны целые положительные числа $A$ и $B$ ($A > B$). На отрезке длины $A$ размещено максимально возможное количество отрезков длины $B$ (без наложений). Используя операцию деления нацело, найти количество отрезков $B$, размещенных на отрезке $A$.
Даны целые положительные числа $A$ и $B$ ($A > B$). На отрезке длины $A$ размещено максимально возможное количество отрезков длины $B$ (без наложений). Используя операцию взятия остатка от деления нацело, найти длину незанятой части отрезка $A$.
Дано двузначное число. Вывести вначале его левую цифру (десятки), а затем — его правую цифру (единицы). Для нахождения десятков использовать операцию деления нацело, для нахождения единиц — операцию взятия остатка от деления.
Дано двузначное число. Найти сумму и произведение его цифр.
Дано двузначное число. Вывести число, полученное при перестановке цифр исходного числа.
Дано трехзначное число. Используя одну операцию деления нацело, вывести первую цифру данного числа (сотни).
Дано трехзначное число. Вывести вначале его последнюю цифру (единицы), а затем — его среднюю цифру (десятки).
Дано трехзначное число. Найти сумму и произведение его цифр.
Дано трехзначное число. Вывести число, полученное при прочтении исходного числа справа налево.
Дано трехзначное число. В нем зачеркнули первую слева цифру и приписали ее справа. Вывести полученное число.
Дано трехзначное число. В нем зачеркнули первую справа цифру и приписали ее слева. Вывести полученное число.
Дано трехзначное число. Вывести число, полученное при перестановке цифр сотен и десятков исходного числа (например, 123 перейдет в 213).
Дано трехзначное число. Вывести число, полученное при перестановке цифр десятков и единиц исходного числа (например, 123 перейдет в 132).
Дано целое число, большее 999. Используя одну операцию деления нацело и одну операцию взятия остатка от деления, найти цифру, соответствующую разряду сотен в записи этого числа.
Дано целое число, большее 999. Используя одну операцию деления нацело и одну операцию взятия остатка от деления, найти цифру, соответствующую разряду тысяч в записи этого числа.
С начала суток прошло $N$ секунд ($N$ — целое). Найти количество полных минут, прошедших с начала суток.
С начала суток прошло $N$ секунд ($N$ — целое). Найти количество полных часов, прошедших с начала суток.
С начала суток прошло $N$ секунд ($N$ — целое). Найти количество секунд, прошедших с начала последней минуты.
С начала суток прошло $N$ секунд ($N$ — целое). Найти количество секунд, прошедших с начала последнего часа.
С начала суток прошло $N$ секунд ($N$ — целое). Найти количество полных минут, прошедших с начала последнего часа.
Дано количество миллиметров $A$ и сантиметров $B$. Вывести число, равное сумме $A + B$ в миллиметрах.
Дано количество граммов $A$ и килограммов $B$. Вывести число, равное сумме $A + B$ в граммах.
Дано четырехзначное число. Найти сумму и произведение его цифр.
Дано четырехзначное число. Вывести число, полученное при прочтении исходного числа справа налево.
Дано четырехзначное число. В нем зачеркнули первую слева цифру и приписали ее справа. Вывести полученное число.
Дано четырехзначное число. В нем зачеркнули первую справа цифру и приписали ее слева. Вывести полученное число.
Дано четырехзначное число. Вывести число, полученное при перестановке первой и последней цифр исходного числа (например, 1234 перейдет в 4231).
